\chapter{総合考察}
本章では，第3章の調査と第5章の検証実験の結果を踏まえ，本研究で得られた知見を総合的に考察する．\todo{検証実験の結果が出たら6.2節・6.3節を中心に加筆する}

\section{推し活の便益と負担は同じ行動から生じる}
第3章の調査では，推し活の頻度と没入が，幸福度と推し疲れの両方と関連していた．また，幸福度と推し疲れは正の相関を示し，推し活によって幸福を感じている人ほど推し疲れも感じやすい傾向があった．このことは，推し活の支援において「活動を減らす」ことを目標にすると，負担だけでなく楽しさや生きがいまで失わせてしまう可能性が高いことを示している．

この知見は，一般的なデジタルウェルビーイングの支援が前提とする「減らすほど良い」という考え方とは異なる．推し活のように本人の生きがいと結びついた活動では，活動の量ではなく，負担につながる特定の関わり方（深夜のSNS閲覧，予算を決めない購入，すべてを追おうとする義務感，他のファンとの比較など）を整えることが重要である．OshiFitのプランが，活動をやめることではなく，時間帯や予算，追う範囲，見る言葉を整える内容になっているのはこのためである．

また，気分の回復のための推し活は幸福度とのみ関連し，推し疲れとは関連しなかった．つらいときに推しに支えられるという推し活の役割は，負担を伴わずに幸福を支える側面であると考えられる．寝る前に推しの曲だけを聴く時間をつくるプランなど，推し活の回復的な側面を生かすプランを含めたのはこのためである．

\section{性格特性の役割}
性格特性は推し活の関わり方と関連していたが，推し活の頻度の効果を調整してはいなかった．また，性格のみでは明確なタイプに分かれなかった．これらの結果は，性格を固定的なタイプとして扱い，タイプごとに決まった推し活を勧めるという方法が，推し活の支援には適さないことを示している．

一方で，性格は推し活の関わり方の違いを通じて，どのような負担が生じやすいかに関わっている可能性がある．例えば，神経症傾向の高い人は気分の回復のために推し活をしやすく，SNS上での比較による落ち込みを経験しやすいと考えられる．OshiFitでは，本人の負担の状況を主な手がかりとしてプランを選び，性格は同じ優先度のプランの中から本人に合いやすいものを選ぶ補助として用いた．この方法は，性格による個人化の利点を生かしつつ，性格による決めつけを避けるものである．

\todo{検証実験後に追記：性格によってプランの実行率や効果に違いがあったか}

\section{不確かさを示すことの意味}
個人の幸福度と推し疲れの予測精度は限られており，特に推し疲れの交差検証決定係数は.11にとどまった．この結果は，個人に見込みを示すシステムにとって重要な制約である．OshiFitでは，見込みを80\%の幅とともに示し，差が小さいときは「ほぼ同じ」と表示し，見込みが診断ではないことを明記した．

不確かさを示すことは，システムの説得力を弱めるように見えるかもしれない．しかし，推し活のように本人の価値観が関わる領域では，不確かな予測を確かなもののように示すことは，本人の判断を誤らせるだけでなく，システムへの信頼を損なうおそれがある．予測の幅を示し，算出式を開示し，最終的な選択を本人に委ねることで，利用者はシステムの提案を一つの材料として自ら考えることができる．

\todo{検証実験後に追記：利用者は幅のある見込みをどのように受け取ったか（自由記述など）}

\section{アバターで2つの未来を示すことについて}
OshiFitは，本人に似たアバターでプランAとプランBの見込みを並べて示した．先行研究\cite{fox2009,hershfield2011}が対象とした運動や貯蓄とは異なり，推し活には一方向の理想がない．そのためOshiFitでは，アバターに一方の未来だけを示すのではなく，今のままの未来と整えた未来を並べ，どちらも選べる形にした．また，アバターの表情は数値の区分だけで決まるようにし，大げさな表現を避けた．

\todo{検証実験後に追記：アバターの提示が選択や実行にどう関わったか．なお本研究の設計では，アバターの有無による比較は行っていないため，アバターそのものの効果は検証の対象外である}

\section{本研究の限界}
本研究には次の限界がある．

\begin{enumerate}
\item \textbf{調査対象の偏り}：第3章の調査の回答者は大学の連絡網を通じて募集した若年層が中心であり，推し活を行う人全体を代表しているとは限らない．年齢や推しのジャンルによって結果が異なる可能性がある．
\item \textbf{測定の信頼性}：TIPI-Jを原版の7件法ではなく5件法で用いたこと，また一部の合成得点（推し活の頻度など）の信頼性係数が低かったことにより，関連の大きさが過小に推定されている可能性がある．
\item \textbf{横断的なデータ}：第3章の調査は横断的であり，示された関連は因果関係を意味しない．
\item \textbf{無作為化を行わない設計}：検証実験の群は本人の状態と選択で決まるため，群の差をプランやOshiFitの効果と断定することはできない．傾向スコアと開始時の値で調整しても，測定されていない違いは残る．
\item \textbf{見込みの仮定の暫定性}：プランBの見込みに用いた想定改善量，複数の提案の倍率（0.8・0.6），提案の選び方の係数（12・18），群Aの基準（60点・5軸以上），個人補正のばらつきの値は，データから推定した値ではなく暫定的な設計上の値であり，実験の結果によって確かめ，更新する必要がある．
\item \textbf{自己報告}：幸福度と推し疲れはいずれも自己報告に基づくため，望ましく見せようとする回答の影響を受ける可能性がある．前研究のような生理指標との組み合わせは，本研究では行っていない．
\item \textbf{期間}：約28日間は，行動が習慣として定着するには短い\cite{lally2010}．当初計画の6週間から短くしたため，提案の効果が現れる前に調査が終わる可能性がある．
\item \textbf{計画からの変更}：当初の研究計画では，無作為の3群（記録のみ，文章での提案，提案とアバター）と2週間の記録だけの期間を予定していたが，D1に基づき本人の選択による群とし，記録だけの期間も設けなかった．そのため，アバターの有無による比較や，2週目の値を共変量とする分析はできない．
\item \textbf{回答の負担と欠測}：毎日のアンケートは負担が小さいように設計したが，回答しない日が生じる．混合モデルは欠けた日があっても推定できるが，回答しなかった理由が状態の悪さと関係する場合には偏りが残る．
\item \textbf{生成AIによる提案文}：提案文は利用者の推しに合わせて生成AIが書くため，同じ型でも文の具体性や分かりやすさが人によって異なる．提案ごとの分析は型の単位で行うが，その結果には文の違いの影響も含まれる．また，生成された文は研究者が事前に一つずつ確認したものではない．禁止した表現の確認と固定の文への切り替えで対処したが，すべての不適切な文を防げるとは限らない．
\item \textbf{デバッグ協力者の参加}：システムのデバッグに協力した研究室の院生が本実験に参加した場合，提案の内容や研究の狙いを事前に知っている．本分析から除いた結果も示すが，完全には取り除けない．
\end{enumerate}

\section{社会的な意義}
本研究は，推し活という多くの人が日常的に行っている活動を対象に，それをやめさせるのではなく，本人が楽しみを保ちながら負担を減らせるよう支援する方法を示した．推し活に限らず，ゲーム，スポーツ観戦，創作活動など，楽しさと負担が表裏一体となった趣味や娯楽は多い．本研究の設計指針（一律に減らさない，行動と性格を組み合わせて個人化する，便益と負担を同時に示す，不確かさを正直に示す）は，これらの活動を持続可能な形に整えるための支援にも応用できると考えられる．
